\clearpage

\AddToShipoutPictureBG{
    \begin{tikzpicture}[remember picture, overlay]
        \node[opacity=0.06] at (current page.center) {
            \includegraphics[width=\paperwidth,height=\paperheight]{image/plante}
        };
    \end{tikzpicture}
}

\chapter{Espace modules}
\label{chap:espace_modules}

\bigskip

\begin{tcolorbox}[colback=blue!6, colframe=blue!55!black, boxrule=0.5pt, arc=4pt,
  left=10pt, right=10pt, top=8pt, bottom=8pt,
  title={\textbf{\small Positionnement du chapitre dans la démonstration}},
  fonttitle=\small\bfseries, coltitle=white, attach boxed title to top left={yshift=-2mm, xshift=4mm},
  boxed title style={colback=blue!55!black, arc=3pt, boxrule=0pt}]
\small
Ce chapitre a une fonction de \textbf{cartographie des apprentissages}. Il rassemble les matériaux de formation (cours, lectures, productions intermédiaires) et montre comment ils ont été mobilisés au fil des phases du Design Thinking.

Il ne constitue pas, à lui seul, la démonstration probante des compétences retenues. La \textbf{démonstration détaillée par traces et preuves} est menée dans le chapitre \emph{Situations professionnalisantes et analyse des compétences}, où les acquis A2, B4 et C2 sont explicitement argumentés à partir de situations, d'artefacts et d'analyses réflexives.
\end{tcolorbox}

\bigskip

La présente section rassemble une sélection représentative des apprentissages dispensés au fil de la formation Innokick. Loin de prétendre à l'exhaustivité, elle réunit un échantillon de traces témoignant d'un engagement actif dans les modules suivis et d'une volonté délibérée d'inscrire ces savoirs dans une réflexion personnelle et une mise en perspective professionnelle.

% ------------------------------------------------------------------------------
\section{Semences théoriques}
% ------------------------------------------------------------------------------

Les notes prises durant les sessions, en présentiel ou en ligne, constituent les apports initiaux de chaque apprentissage. S'y ajoutent les résumés personnels élaborés \textit{a posteriori} ainsi que les documents fournis par les intervenants, qui ont servi de matériau de base à la compréhension.

% ------------------------------------------------------------------------------
\section{Ramifications conceptuelles}
% ------------------------------------------------------------------------------

Les mindmaps et schémas réalisés au cours de la formation rendent visibles les liens établis entre les concepts. Ces représentations visuelles attestent d'un effort de synthèse et d'appropriation qui dépasse la simple réception passive des contenus.

% ------------------------------------------------------------------------------
\section{Lectures nourrici\`ieres}
% ------------------------------------------------------------------------------

Les synthèses de lectures complémentaires, effectuées dans le cadre des quinze cours choisis hors de mon domaine d'apprentissage habituel, ont constitué l'apport intellectuel nécessaire au développement de compétences jusqu'alors peu cultivées. Chaque lecture a élargi le champ de références mobilisables en dehors du registre technique.

% ------------------------------------------------------------------------------
\section{Productions et récoltes intermédiaires}
% ------------------------------------------------------------------------------

Les exercices, livrables de groupe et productions individuelles réalisés durant les modules constituent les premiers résultats tangibles de cette démarche. Plus ou moins aboutis selon les modules, ils jalonnent la progression et offrent des preuves observables de l'avancement.

% ------------------------------------------------------------------------------
\section{Boutures réflexives}
% ------------------------------------------------------------------------------

Enfin, les annotations, questionnements et connexions établies entre les contenus théoriques et ma pratique professionnelle chez TWIICE constituent autant de transpositions : des fragments de réflexion déplacés d'un contexte académique vers un terrain professionnel, afin d'en éprouver la pertinence dans la réalité du travail.

% ------------------------------------------------------------------------------
\section{Organisation}
% ------------------------------------------------------------------------------

L'ensemble de ces éléments est ordonné selon les phases effectivement traversées durant la période couverte par ce portfolio. Chaque élément est, dans la mesure du possible, relié à une ou plusieurs compétences du référentiel Innokick, afin de faciliter la démonstration de la mobilisation des ressources lors de l'examen.

% ==============================================================================
\section{Phase 1 : Empathie}
% ==============================================================================

\competences{A1 (Curiosité et esprit critique), B1 (Recherche et analyse), C1 (Interaction parties prenantes)}

\subsection{Semences théoriques}

Les cours d'empathie utilisateur et d'observation terrain m'ont fourni une méthode explicite pour distinguer ce qui est dit par l'utilisateur et ce qui est interprété par le concepteur. L'apport clé de cette phase a été de ralentir volontairement l'entrée en solution afin de qualifier d'abord le besoin.

\subsection{Lectures nourricières}

Deux familles de lectures ont soutenu cette phase :
\begin{itemize}
  \item des textes sur l'écoute active et la reformulation des besoins, utiles pour la compétence C1 ;
  \item des ressources sur l'observation contextuelle, utiles pour passer d'une intuition technique à une hypothèse argumentée (A1, B1).
\end{itemize}

\subsection{Productions et récoltes intermédiaires}

Les productions de la phase Empathie sont principalement des notes d'entretien, des synthèses de signaux faibles et des tableaux de besoins reformulés. Leur utilité est d'avoir servi de matière première aux phases suivantes, en particulier à la formulation des problématiques en phase Définition.

% ==============================================================================
\section{Phase 2 : Définition}
% ==============================================================================

\competences{A2 (Gestion des divergences), B2 (Méthodes de formulation), C2 (Synthèse et opportunités)}

% ------------------------------------------------------------------------------
\subsection{Semences théoriques}
% ------------------------------------------------------------------------------

\paragraph{Cours de méthodes agiles et gestion de projet}
Les mécanismes de diver\-gence-con\-vergence appliqués à la prise de décision collective : timeboxing, \textit{definition of done}, priorisation par critères explicites. Ce cours a fourni les bases d'une approche structurée pour transformer un débat d'opinions en processus de décision traçable.

\paragraph{Cours de communication et médiation}
Les techniques de reformulation et de recadrage en situation de désaccord. Apprentissage clé : recentrer un groupe sur l'objectif commun plutôt que sur les positions individuelles.

\paragraph{Cours de méthodes de recherche qualitative et quantitative}
La construction d'un protocole de recherche minimal (question, échantillon, données, critères) avant toute conclusion. Ce socle méthodologique a permis de comprendre la différence entre chercher une solution et comprendre un problème.

\paragraph{Cours de pensée systémique}
La capacité à voir au-delà du symptôme immédiat et à relier un choix isolé à ses conséquences en cascade (qualité, traçabilité, charge documentaire, calendrier).

\paragraph{Cours de veille et intelligence économique}
Les méthodes de veille structurée (sources, fréquence, outils, diffusion) pour transformer une recherche ponctuelle en démarche systématique et reproductible.

% ------------------------------------------------------------------------------
\subsection{Lectures nourric\`ieres}
% ------------------------------------------------------------------------------

\begin{itemize}
  \item \textbf{\citet{adams2004}}, \textit{Change Your Questions, Change Your Life} : le Question Thinking : transformer les questions de jugement en questions d'apprentissage.
  \item \textbf{\citet{jakobiak1998}}, \textit{L'intelligence économique en pratique} : la veille informationnelle comme processus continu et proactif.
  \item \textbf{\citet{senge1990}}, \textit{The Fifth Discipline} : la pensée systémique appliquée à l'analyse des interdépendances organisationnelles.
  \item \textbf{\citet{brown2009}}, \textit{Change by Design} : la pratique du Design Thinking comme posture d'empathie outillée, ancrage théorique des phases 1 et 2 du présent chapitre.
\end{itemize}

% ------------------------------------------------------------------------------
\subsection{Ramifications conceptuelles}
% ------------------------------------------------------------------------------

\subsubsection*{Matrice d'Eisenhower}

Outil de tri à deux axes (importance / urgence) présenté en cours pour prioriser des décisions concurrentes. Son application concrète au cas TWIICE (certification IEC~62304) est documentée plus bas comme \textbf{Illustration~1} dans la section \textit{Productions et récoltes intermédiaires}, afin d'éviter de présenter deux fois le même tableau.

\subsubsection*{Grille \textit{build vs buy}}

Appliquée aux sujets identifiés comme prioritaires par la matrice d'Eisenhower, avec des critères explicites : complexité technique, coût, intégration, traçabilité, délai de certification, dépendance fournisseur.

\subsubsection*{Méthode des 5 Pourquoi (\textit{5~Why})}

Appliquée rétrospectivement à la situation de la banque privée pour remonter à la cause racine du surcoût des rapports :

\begin{enumerate}
  \item Pourquoi les rapports coûtent-ils cher ? $\rightarrow$ Personnalisation manuelle.
  \item Pourquoi manuelle ? $\rightarrow$ Pas de template standardisé.
  \item Pourquoi pas de template ? $\rightarrow$ Exigences clients non catégorisées.
  \item Pourquoi non catégorisées ? $\rightarrow$ Aucune analyse des points communs.
  \item Pourquoi pas d'analyse ? $\rightarrow$ Processus centré sur la réponse individuelle.
\end{enumerate}

% ------------------------------------------------------------------------------
\subsection{Productions et récoltes intermédiaires}
% ------------------------------------------------------------------------------

La situation retenue pour illustrer concrètement la mobilisation des acquis de la phase Définition est celle de la \textbf{priorisation des décisions logicielles chez TWIICE}, dans le cadre de la certification IEC~62304 de l'exosquelette (compétence C2). Ce cas concentre à lui seul trois outils enseignés en cours, transposés directement dans le contexte professionnel.

\subsubsection{Illustration 1 : Matrice d'Eisenhower formalisée}

L'équipe logicielle faisait face à une accumulation de décisions techniques parallèles, sans hiérarchie. La matrice d'Eisenhower, découverte lors du cours de gestion de l'innovation et priorisation, a été appliquée pour trier ces décisions :

\begin{table}[h!]
  \centering
  \begin{tabular}{|p{3cm}|p{5.5cm}|p{5.5cm}|}
    \hline
    & \textbf{Urgent} & \textbf{Non urgent} \\
    \hline
    \textbf{Important}
      & Traçabilité des SOUP et artefacts logiciels : bloque la certification IEC~62304
      & Standardisation des templates de documentation : gain de productivité à moyen terme \\
    \hline
    \textbf{Non important}
      & Choix microcontrôleur vs Jetson pour l'interface utilisateur : débat actif mais non bloquant
      & Évaluation d'outils DevOps\footnote{Le \textbf{DevOps} (contraction de \emph{Development} et \emph{Operations}) désigne un ensemble de pratiques visant à rapprocher les équipes de développement et d'exploitation, et à automatiser le cycle complet de livraison du logiciel (construction, tests, déploiement, supervision). Les \emph{outils DevOps} regroupent les plateformes qui supportent cette chaîne : gestion de code source, intégration continue, conteneurisation, infrastructure-as-code, supervision.} complémentaires : confort d'équipe, aucun impact certification \\
    \hline
  \end{tabular}
  \caption{Matrice d'Eisenhower formalisée : Décisions logicielles TWIICE}
  \label{tab:eisenhower-illustration1}
\end{table}

\textbf{Résultat} : l'équipe a immédiatement concentré son énergie sur le quadrant urgent + important et planifié les sujets restants dans les sprints suivants.

\subsubsection{Illustration 2 : Grille \textit{build vs buy}}

Une fois la traçabilité SOUP identifiée comme priorité, la question concrète était : construire un système de stockage et de versionnage\footnote{Le \textbf{versionnage} (ou \emph{versioning}) désigne la gestion systématique des différentes versions successives d'un fichier ou d'un logiciel. Chaque modification est enregistrée avec un identifiant unique, ce qui permet de retrouver l'état exact du code à une date donnée, de comparer des versions, ou de revenir en arrière. C'est une exigence fondamentale de la traçabilité réglementaire en contexte médical.\label{fn:versionnage}} sur mesure, ou s'appuyer sur des solutions existantes (GitLab\footnote{\textbf{GitLab} est une plateforme web intégrée de gestion de code source (basée sur Git) et d'industrialisation du développement logiciel : suivi des tickets, revues de code, intégration continue, registre d'artefacts, gestion des déploiements. Elle existe en version \emph{cloud} (hébergée par l'éditeur) ou \emph{self-hosted} (hébergée sur les serveurs internes de l'organisation).\label{fn:gitlab}}, JFrog\footnote{\textbf{JFrog Artifactory} est un gestionnaire d'artefacts logiciels : un entrepôt centralisé qui stocke, versionne et trace les binaires produits par le développement (bibliothèques, paquets, images de conteneurs), avec des métadonnées et des contrôles d'intégrité. Il répond à une exigence forte de la norme IEC~62304 sur la maîtrise des composants logiciels déployés.\label{fn:jfrog}}) ? La grille suivante a structuré la comparaison :

\begin{table}[h!]
  \centering\small
  \begin{tabularx}{\textwidth}{|p{3.6cm}|X|X|X|}
    \hline
    \textbf{Critère} & \textbf{GitLab (cloud)} & \textbf{Solution custom} & \textbf{JFrog Artifactory} \\
    \hline
    \textbf{Complexité d'intégration}
      & Faible : déjà utilisé par l'équipe
      & Élevée : développement et maintenance internes
      & Moyenne : intégration CI/CD\footnote{\textbf{CI/CD} désigne l'\emph{intégration continue} (\emph{Continuous Integration}) et la \emph{livraison/déploiement continu(e)} (\emph{Continuous Delivery / Deployment}) : une chaîne automatisée dans laquelle chaque modification de code déclenche la construction, les tests et, le cas échéant, le déploiement du logiciel. Elle garantit que le code reste en permanence dans un état livrable et traçable.\label{fn:cicd}} à configurer \\
    \hline
    \textbf{Coût}
      & Inclus dans l'abonnement existant
      & Élevé (heures de développement internes)
      & Licence supplémentaire à négocier \\
    \hline
    \textbf{Traçabilité IEC~62304}
      & Bonne : versionnage natif, tags, releases
      & Totale : contrôle complet du format
      & Bonne : métadonnées et checksums natifs \\
    \hline
    \textbf{Délai de mise en \oe{}uvre}
      & Immédiat
      & Plusieurs semaines/mois
      & Quelques jours \\
    \hline
    \textbf{Dépendance fournisseur}
      & Moyenne : données exportables, mais cloud
      & Nulle : 100\,\% interne
      & Élevée : solution propriétaire \\
    \hline
    \textbf{Risque cybersécurité (dispositif médical)}
      & Modéré : hébergement cloud, conformité à vérifier
      & Faible : hébergement interne
      & Modéré : options self-hosted disponibles \\
    \hline
  \end{tabularx}
  \caption{Grille build vs buy : Traçabilité SOUP (IEC~62304)}
  \label{tab:buildvsbuy}
\end{table}

\textbf{Résultat} : cette grille a permis de dépasser le débat d'opinions et de converger vers une décision fondée sur des critères explicites et pondérés, alignés avec les exigences de certification.

\subsubsection{Preuves complémentaires (en annexe)}

Les documents suivants, trop volumineux pour être intégrés dans le corps du portfolio, sont reproduits en annexe. Ils sont mobilisés comme \emph{preuves} au sens de Tardif dans le chapitre \textit{Situations professionnalisantes} ; ils ne sont listés ici qu'à titre de carte des artefacts disponibles :

\begin{itemize}
  \item \textbf{Annexe~\ref{annex:ndt-sw-001}} : note de décision technique ; formalise le choix retenu, les options écartées et les risques résiduels (exigence IEC~62304).
  \item \textbf{Annexe~\ref{annex:backlog}} : backlog priorisé ; montre la hiérarchisation des tâches après application de la matrice d'Eisenhower.
\end{itemize}

% ==============================================================================
\section{Phase 3 : Idéation}
% ==============================================================================

\competences{A3 (Posture créative), B3 (Techniques d'idéation), C3 (Évaluation des idées)}

\subsection{Semences théoriques}

La phase Idéation a consolidé un changement de posture : explorer simultanément plusieurs pistes imparfaites plutôt qu'une seule solution précoce, et conserver ces variantes en parallèle avant de retenir les plus prometteuses. Les outils vus en cours (brief design, divergence-convergence, tri collectif) ont rendu ce changement opératoire.

\subsection{Ramifications conceptuelles}

Trois principes ont guidé l'idéation :
\begin{itemize}
  \item suspendre temporairement le jugement pendant la divergence ;
  \item expliciter les critères de tri avant de converger ;
  \item tester la qualité d'une idée sur sa valeur d'usage, pas uniquement sur son èlegance technique.
\end{itemize}

\subsection{Productions et récoltes intermédiaires}

Les récoltes de cette phase sont constituées de variantes de concepts, de versions successives de pitch et de traces de sélection d'idées. Même lorsque certaines pistes ont été élaguées, elles ont joué un rôle utile : rendre explicite le raisonnement ayant conduit à l'option retenue.

% ==============================================================================
\section{Phase 4 : Prototypage}
% ==============================================================================

\competences{A4 (Agilité et adaptabilité), B4 (Méthodes agiles et lean), C4 (Évaluation du projet)}

% ------------------------------------------------------------------------------
\subsection{Semences théoriques}
% ------------------------------------------------------------------------------

\paragraph{Cours de méthodes agiles et prototypage itératif}
Introduction aux cycles sprint, revue et rétrospective, ainsi qu'à la logique lean de réduction du gaspillage dans le processus de développement. Graine centrale plântée ici : la valeur d'un livrable ne se mesure pas à sa complétude, mais à ce qu'il permet d'apprendre sur les hypothèses sous-jacentes.

Les enjeux de traçabilité IEC~62304 et d'articulation entre agilité et certification sont traités dans cette phase comme des apprentissages issus de l'expérience terrain chez TWIICE, et non comme des contenus de cours.


% ------------------------------------------------------------------------------
\subsection{Lectures nourric\`ieres}
% ------------------------------------------------------------------------------

\begin{itemize}
  \item \textbf{\citet{beck2001}}, \textit{Manifesto for Agile Software Development} : les quatre valeurs fondatrices de l'agilité (individus et interactions, logiciel opérationnel, collaboration client, adaptation au changement). Point de référence axiologique mobilisé en arrière-plan de toute la phase~4.
  \item \textbf{\citet{sutherland2014}}, \textit{Scrum\footnote{\textbf{Scrum} est un cadre de travail (\emph{framework}) agile organisant le développement par cycles courts appelés \emph{sprints} (typiquement 2 à 4 semaines). Il définit trois rôles (Product Owner, Scrum Master, Équipe de développement), quatre événements (Sprint Planning, Daily, Review, Rétrospective) et trois artefacts (Product Backlog, Sprint Backlog, Increment). Son objectif est de livrer fréquemment de la valeur tout en réduisant le risque par l'inspection et l'adaptation continues.\label{fn:scrum}} : The Art of Doing Twice the Work in Half the Time} : la philosophie fondatrice du sprint comme unité d'apprentissage, pas seulement de livraison.
  \item \textbf{\citet{poppendieck2003}}, \textit{Lean Software Development} : sept principes lean appliqués au logiciel, dont l'élimination des gaspillages documentaires inutiles et la décision au dernier moment responsable.
  \item \textbf{\citet{humble2010}}, \textit{Continuous Delivery} : la pipeline\footnote{Une \textbf{pipeline} (litt. \emph{pipeline} = conduite) désigne, en développement logiciel, une chaîne d'étapes automatisées déclenchées à chaque modification du code source : compilation, tests automatisés, analyse de qualité, packaging, déploiement. Elle matérialise concrètement la démarche CI/CD.} CI/CD comme infrastructure de récolte continue : chaque commit\footnote{Un \textbf{commit} désigne l'enregistrement d'un ensemble de modifications de code source dans le système de versionnage, accompagné d'un message décrivant l'intention. C'est l'unité élémentaire de traçabilité du développement.} produit un artefact vérifié, traçable, reproductible.
  \item \textbf{\citet{mchugh2014}}, \emph{An Agile Implementation within a Medical Device Software Organisation} : étude de cas appliquée d'une transition agile sous contrainte IEC~62304, point d'appui direct pour la situation TWIICE documentée au chapitre suivant.
\end{itemize}

% ------------------------------------------------------------------------------
\subsection{Ramifications conceptuelles}
% ------------------------------------------------------------------------------

\subsubsection*{Scrum adapté au développement médical (compétence B4)}

La contrainte IEC~62304 oblige à reconfigurer le cycle agile standard. L'adaptation réalisée chez TWIICE consiste à intégrer les exigences documentaires au rythme sprint, sans en briser la cadence :

\begin{table}[h!]
  \centering\small
  \begin{tabularx}{\textwidth}{|p{3.2cm}|X|X|}
    \hline
    \textbf{Événement Scrum} & \textbf{Adaptation standard} & \textbf{Adaptation médicale (IEC~62304)} \\
    \hline
    Sprint Planning
      & Sélectionner les items du backlog selon la valeur utilisateur
      & Ajouter un critère : risque de certification (bloquer la validation = priorité haute) \\
    \hline
    Sprint Review
      & Démontrer le logiciel fonctionnel
      & Vérifier que les enregistrements qualité associés sont produits et archivés \\
    \hline
    Rétrospective
      & Améliorer le processus d'équipe
      & Identifier les dettes documentaires accumulées et planifier leur remboursement \\
    \hline
    Definition of Done
      & Code + tests passants
      & Code + tests + documentation IEC~62304 + revue qualité + archivage GitLab \\
    \hline
  \end{tabularx}
  \caption{Scrum adapté au contexte de certification médicale IEC~62304 : TWIICE}
  \label{tab:scrum-medical}
\end{table}

\subsubsection*{Priorisation du backlog : valeur $\times$ risque certification}

La priorisation classique par valeur utilisateur ne suffit pas dans un contexte réglementé. Une fonction à forte valeur mais à risque de certification élevé doit monter dans le backlog pour éviter les surprises tardives. La formule appliquée :

\begin{quote}
  \textit{Priorité} = Valeur fonctionnelle $\times$ Risque certification
\end{quote}

Un item de faible valeur fonctionnelle mais dont la validation par l'auditeur est imposée en phase précoce est traité avant un item plus désirable mais moins urgent pour la certification.

% ------------------------------------------------------------------------------
\subsection{Productions et récoltes intermédiaires}
% ------------------------------------------------------------------------------

La récolte principale de la phase Prototypage, pour la compétence B4, est documentée dans le chapitre \textit{Situations professionnalisantes} (section B4) et dans les annexes du portfolio. Elle comprend trois pièces cueillies directement du terrain professionnel chez TWIICE :

\begin{itemize}
  \item \textbf{Annexe~\ref{annex:backlog} : Backlog priorisé} ; extrait du board GitLab après application de la grille de priorisation fondée sur la valeur et le risque. Cette récolte documente le passage d'un catalogue de débats à une file de travail ordonnée et actionnable.
  \item \textbf{Annexe~\ref{annex:sdp-markdown} : Software Development Plan} ; document-cadre qui formalise la structure sprint-revue-rétrospective adaptée à l'IEC~62304, ainsi que les règles de traçabilité et de gouvernance associées.
  \item \textbf{Annexe~\ref{annex:retro-b4} : Compte-rendu de la première rétrospective} ; preuve de la boucle d'amélioration continue, avec actions Start/Stop/Continue, responsables et critères de vérification.
\end{itemize}

% ==============================================================================
\section{Phase 5 : Test}
% ==============================================================================

\competences{A5 (Évaluation des performances), B5 (Protocoles de test), C5 (Analyse des résultats)}

\subsection{Semences théoriques}

La phase Test a renforcé la distinction entre vérifier que "\c ca fonctionne" et vérifier que "\c ca répond à l'exigence". Les cours ont insisté sur la nécessité de définir les critères d'acceptation avant l'exécution des tests.

\subsection{Protocoles et limites}

Les protocoles mobilisés couvrent trois niveaux :
\begin{itemize}
  \item tests unitaires pour valider les comportements locaux ;
  \item tests d'intégration pour contrôler les interfaces ;
  \item tests système pour confirmer la cohérence globale avec les exigences.
\end{itemize}

La limite principale identifiée concerne la couverture en contexte réel : certains scénarios restent sensibles à l'environnement d'usage et nécessitent des itérations complémentaires.

\subsection{Productions et récoltes intermédiaires}

Les traces de cette phase incluent des plans de test, des critères de validation, des résultats d'exécution et des notes de non-conformité. Ces artefacts ont servi à prioriser les corrections et à objectiver les arbitrages avant déploiement.

% ==============================================================================
\section{Phase 6 : Implémentation}
% ==============================================================================

\competences{A6 (Communication et promotion), B6 (Planification de mise en \oe{}uvre), C6 (Scénarios d'implémentation)}

Cette phase est \textbf{hors du périmètre temporel d'évaluation} couvert par le module MDM. Elle correspond à un horizon de déploiement (mise en service, communication produit, scénarios d'industrialisation) qui se situe \emph{après} la période documentée par ce portfolio et qui relève des modules ultérieurs du cursus Innokick. Le Plan d'Études Cadre du Master HES-SO Innokick (juillet~2024, annexe~\ref{annex:pec-innokick}) n'impose pas la couverture exhaustive des six phases du Design Thinking au titre du seul module MDM : le PEC articule la maturation des compétences au fil des modules successifs, la mise en \oe{}uvre opérationnelle étant explicitement adossée aux modules d'aval. Le présent portfolio se conforme donc à ce découpage modulaire.

Ce choix de périmètre est assumé : les preuves mobilisables dans le cadre du présent portfolio concernent les phases~1 à~5, qui sont les seules à avoir été \emph{réellement traversées} sur la période d'évaluation. Documenter la phase~6 sans l'avoir vécue aurait conduit à produire un discours sur l'intention plutôt qu'une trace au sens de Tardif --- ce que le cadre méthodologique adopté exclut explicitement.

Le chapitre \textit{Perspectives pour la suite} reprend cette phase sous l'angle prospectif (passage vers la \textit{business layer}, certification PRINCE2 comme appui) sans prétendre en faire la démonstration probatoire ici.

\section*{Transition vers le chapitre de démonstration}

La démonstration probante des compétences A2, B4 et C2, appuyée sur des situations, des traces et des preuves détaillées, est présentée dans le chapitre suivant, \textit{Situations professionnalisantes et analyse des compétences}.

\clearpage
\ClearShipoutPictureBG

